\documentclass[../main]{subfiles}
\begin{document}
\chapter{Mantenimiento}
\img{images/07/mantenimiento.jpg}{Mantenimiento predictivo: medición
de vibraciones en acople rígido.}
\info{Mantenimiento}{Conservación de una cosa en buen estado o en una
situación determinada para evitar su degradación.}
\section{Tipos de fallas}
\info{Falla}{Es la incapacidad del activo de cumplir con su función,
cualquier función. Son todos los estados indeseables del sistema o equipo}
\subsection{Falla técnica}
Una falla técnica a diferencia de la falla funcional permite que la
máquina siga funcionando, pero con algunas alteraciones que son
perceptibles, ruido, vibraciones, y que en definitiva van a producir
un fallo funcional en algún momento.
\subsection{Falla funcional}
Deja de cumplir con la función, sin embargo no necesariamente
significa que el activo deja de producir. Aquí es muy importante
definir los requerimientos de desempeño. Son resultado de la
incapacidad de los equipos para dar su nivel de desempeño esperado.
Son un síntoma o evento observable por el personal de producción y
mantenimiento.
\section{Tipos de mantenimiento}
\subsection{Mantenimiento correctivo}
Se denomina mantenimiento correctivo a aquel que corrige los defectos
observados en los equipamientos o instalaciones, es la forma más
básica de mantenimiento y consiste en localizar averías o defectos
para corregirlos o repararlos. Cuando hay un fallo funcional se
repara inmediatamente la máquina.
\subsection{Mantenimiento preventivo}
En las operaciones de mantenimiento el mantenimiento preventivo es el
destinado a la conservación de equipos o instalaciones mediante la
realización de revisión y limpieza que garanticen su buen
funcionamiento y fiabilidad. El mantenimiento preventivo se realiza
en equipos en condiciones de funcionamiento, por oposición al
mantenimiento correctivo que repara o pone en condiciones de
funcionamiento aquellos que dejaron de funcionar o están dañados.

\subsection{Mantenimiento predictivo}
Son una serie de acciones que se toman, y técnicas que se aplican,
con el objetivo de detectar posibles fallos y defectos de maquinaria
en las etapas incipientes, para evitar que estos fallos se
manifiesten en uno más grande durante su funcionamiento, evitando que
ocasionen paros de emergencia y tiempos muertos, causando impacto
financiero negativo. Su misión es conservar un nivel de servicio
determinado en los equipos programando las revisiones en el momento
más oportuno. Suele tener un carácter sistemático, es decir, se
interviene aunque el equipo no haya dado ningún síntoma de tener problemas.

\subsection{Mantenimiento legal}

Se denomina mantenimiento legal, al tipo de mantenimiento preventivo
obligatorio que las legislaciones sobre seguridad de equipos e
instalaciones industriales obligan que deben realizarse de forma
periódica por parte de empresas o personal autorizado ajenas a la
empresa propiedad de las instalaciones o equipos.

\subsection{Mantenimiento subcontratado}
También llamado subcontratación o externalización, es un proceso
económico empresarial en el cual una empresa delega ciertas tareas,
no propias de su función pero igualmente importantes e
indispensables, a una empresa externa especializada en la prestación
de esos servicios.

\section{Ficha Técnica de Equipo}

Para una gestión eficiente del mantenimiento, es fundamental contar con una Ficha de Equipo que centralice toda la información técnica y el plan de intervenciones.

\begin{center}
\begin{tabular}{|p{0.45\textwidth}|p{0.45\textwidth}|}
\hline
\multicolumn{2}{|c|}{\textbf{DEPARTAMENTO DE MANTENIMIENTO - FICHA TÉCNICA DE EQUIPO}} \\ \hline
\textbf{Información General} & \textbf{ID del Activo:} [Código único, ej. MOT-001] \\ \hline
\textbf{Nombre del Equipo:} [Ej. Bomba Centrífuga] & \textbf{Criticidad:} [Alta / Media / Baja] \\ \hline
\textbf{Marca:} [Fabricante] & \textbf{Modelo:} [Número de Modelo] \\ \hline
\textbf{Número de Serie:} [S/N] & \textbf{Fecha de Instalación:} [DD/MM/AAAA] \\ \hline
\end{tabular}

\vspace{0.3cm}

\begin{tabular}{|p{0.9\textwidth}|}
\hline
\textbf{Especificaciones Técnicas} \\ \hline
$\bullet$ \textbf{Potencia/Capacidad:} [Ej. 5 HP / 2kW] \\ 
$\bullet$ \textbf{Voltaje/Corriente:} [Ej. 380V / 3 Fases] \\ 
$\bullet$ \textbf{RPM / Velocidad:} [Ej. 1750 RPM] \\ 
$\bullet$ \textbf{Dimensiones:} [Largo x Ancho x Alto] \\ 
$\bullet$ \textbf{Peso:} [kg] \\ 
$\bullet$ \textbf{Tipo de Lubricante:} [Ej. ISO VG 68] \\ \hline
\end{tabular}

\vspace{0.3cm}

\begin{tabular}{|p{0.3\textwidth}|p{0.2\textwidth}|p{0.2\textwidth}|p{0.2\textwidth}|}
\hline
\textbf{Tarea} & \textbf{Frecuencia} & \textbf{Responsable} & \textbf{Ref. Procedimiento} \\ \hline
Inspección Visual & Semanal & Operador & Manual Sec. 2.1 \\ \hline
Lubricación & Mensual & Técnico & Manual Sec. 4.5 \\ \hline
Pruebas Eléctricas & Trimestral & Electricista & Norma IEC 60364 \\ \hline
Overhaul Completo & Anual & Empresa Ext. & Contrato Servicio \\ \hline
\end{tabular}

\vspace{0.3cm}

\begin{tabular}{|p{0.9\textwidth}|}
\hline
\textbf{Contacto del Proveedor/Fabricante} \\ \hline
\textbf{Empresa:} [Nombre] \\ 
\textbf{Persona de Contacto:} [Nombre] \\ 
\textbf{Teléfono/Email:} [Información de contacto] \\ \hline
\end{tabular}

\vspace{0.3cm}

\begin{tabular}{|p{0.9\textwidth}|}
\hline
\textbf{Observaciones:} \\ 
\vspace{1cm} \\ \hline
\end{tabular}
\end{center}

\actividad { Cuestionario: investigar}
\begin{enumerate}
  \item  ¿Qué es la obsolescencia programada?
  \item ¿Qué es la obsolescencia tecnológica?
  \item ¿Para que se utiliza el análisis de criticidad en mantenimiento?
  \item ¿Qué es la burocracia?
  \item ¿Qué es la codificación de equipos?
  \item ¿Qué es el código de barra, QR cómo funcionan y para que
    puede utilizarse?
  \item ¿Qué es el mantenimiento total?

\end{enumerate}

\end{document}